% ============================================================
% PHẦN 3.13 — RỦI RO, GIẢ ĐỊNH VÀ KẾ HOẠCH KIỂM CHỨNG
% ------------------------------------------------------------
% Thuộc: PA3-04 (rủi ro vận hành) + PA3-06 (tổng hợp giả định)
% Người phụ trách chính: Nguyễn Ngọc Cảnh – Operations Lead
% Kiểm tra logic tổng thể: Nguyễn Anh Thư
%
% CÂU HỎI CẦN TRẢ LỜI:
%   - Khi có lỗi video, dữ liệu hoặc câu trả lời AI thì xử lý ra sao?
%   - Khi nào nhóm dừng một chiến dịch hoặc thay đổi thông điệp?
%   - Thị trường mục tiêu có thống nhất với chương Sản phẩm không?
%   - Thông điệp có nói đúng năng lực MVP không?
%   - Timeline có phụ thuộc vào tính năng chưa hoàn thiện không?
%   - Các số liệu là dữ liệu thật hay giả định?
%
% OUTPUT BẮT BUỘC:
%   1. Bảng rủi ro vận hành:
%      Rủi ro | Xác suất | Tác động | Biện pháp phòng ngừa | Người xử lý
%
%   2. Danh sách giả định cần kiểm chứng:
%      (Tất cả thông tin trong Chương 3 được ghi là GIẢ ĐỊNH phải tổng hợp tại đây)
%
%   3. Kế hoạch kiểm chứng từng giả định:
%      Giả định | Cách kiểm chứng | Thời điểm kiểm chứng | Tiêu chí xác nhận
%
% TIÊU CHÍ HOÀN THÀNH:
%   - Mỗi hoạt động phải tạo ra đầu ra kiểm tra được
%   - Phải phân biệt: thông tin xác nhận vs. giả định cần kiểm chứng
%   - Kế hoạch phù hợp với giai đoạn pilot có hỗ trợ
% ============================================================

\subsection{Rủi ro, giả định và kế hoạch kiểm chứng}

% TODO (Nguyễn Ngọc Cảnh): Viết bảng rủi ro vận hành
% TODO (Nguyễn Anh Thư): Tổng hợp toàn bộ giả định từ 3.1–3.12 vào danh sách
% Cả nhóm review chéo để xác nhận giả định vs. thông tin đã kiểm chứng
